\SourceLeaf{249}%
\noindent in abstracto betrachtet widerspruchsfrei ist und durchaus höheren Wesen zukommen könnte. Um aber aus dieser Höhe wieder auf die Erde zurückzukommen, so ist es schade, dass sich die Gleichheit der Volumina räumlicher, bloß symmetrischer, aber nicht kongruenter Gebilde nur mit der Exhaustionsmethode und nicht ebenso elementar beweisen lässt, wie Sie es meines Wissens zuerst für die Fläche des sphärischen Dreiecks gezeigt haben. Ihr Beweis findet sich in etwas anderer Form auch bei \textsc{Lobatschewsky}. Es ist jedoch unklar, ob er Ihren Aufsatz kannte. Jedes sphärische Dreieck lässt sich nämlich in drei gleichschenklige Dreiecke zerlegen. Wendet man diese Zerlegung auf zwei symmetrische Dreiecke an, so ergeben sich in beiden gleiche, nur unterschiedlich angeordnete Dreiecke. (Für die Vollständigkeit müsste man selbstverständlich drei Fälle unterscheiden. Der zweite und der dritte lassen sich auch so ausdrücken, dass je eines der Teildreiecke negativ oder 0 sein kann.)
\mbox{.\ .\ .\ .\ .\ .\ .\ .}

\par\bigskip

\Needspace{7\baselineskip}
\begin{center}
[\textsc{Gerling} \emph{an} \textsc{Gauss}.\quad \emph{Marburg, 15. April 1844.}]
\end{center}

\{Obwohl ich Ihren Brief vom 8. dieses Monats, für den ich Ihnen vorläufig herzlich danke, eigentlich noch gar nicht beantworten kann, weil ich bisher nicht dazu gekommen bin, ihn so zu studieren und aufzunehmen, wie es sein reicher Inhalt und mein Interesse am Gegenstand verlangen, kann ich mir doch das Vergnügen nicht versagen, Ihnen zu einer einzelnen Äußerung darin ein paar Zeilen zu schreiben.

Sie sagen: »Es ist schade, dass sich die Gleichheit der Volumina räumlicher, bloß symmetrischer, aber nicht kongruenter Gebilde nur mit der Exhaustionsmethode und nicht ebenso elementar beweisen lässt«
usw.

Erlauben Sie mir, Ihnen hier einen solchen Beweis mitzuteilen und Sie zu bitten, mir gelegentlich Ihr Urteil darüber zu schreiben.

\textls{Lehrsatz}: Zwei beliebige symmetrische Polyeder haben gleiches Volumen.

\textls{Beweis}: I. Durch geeignet gelegte Ebenen lassen sie sich\space%
